\originalchapter{8.012 进阶作业与项目}\label{ch:advancedps}
本章对应 2008 秋十份作业. 原卷要求独立写解、注明讨论对象、姓名及习题课班号, 禁止提交抄来的答案; 这里作为历史考核要求保留. OCW 作业页没有给出全套作业答案, 以下十三道完整自编题均为 AI 独立解答. 第 10 套为可选不收作业, 原文字称会在线发布答案, 不等于本次已取得该答案.
\originalsection{教材引用题索引}
以下为原卷只给教材题号的 85 题, 不是已经取得的完整题面. 序号逐项对应原卷, 不补造题目或答案. 原卷的额外说明保留在表后.
\noindent PS1: 1: 1.2, 2: 1.7, 3: 1.12, 4: 1.13, 5: 1.14, 6: 1.15, 7: 1.18, 8: 1.20.\par
\noindent PS2: 1: 2.1, 2: 2.3, 3: 2.5, 4: 2.7, 5: 2.9, 6: 2.11, 7: 2.12, 8: 2.15, 9: 2.17, 10: 2.24.\par
\noindent PS3: 1: 2.22, 2: 2.28, 3: 2.30, 4: 2.31, 5: 2.34, 6: 2.35, 7: 3.1, 8: 3.4, 9: 3.5.\par
\noindent PS4: 1: 3.7, 2: 3.11, 3: 3.13, 4: 3.14, 5: 3.15, 6: 3.16, 7: 3.18, 8: 3.20.\par
\noindent PS5: 1: 4.4, 2: 4.5, 3: 4.7, 4: 4.8, 5: 4.9, 6: 4.10, 7: 4.13, 8: 4.20, 9: 4.28.\par
\noindent PS6: 1: 5.4, 2: 5.7, 3: 6.1, 4: 6.2, 5: 6.4, 6: 6.6.\par
\noindent PS7: 1: 6.8, 2: 6.13, 3: 6.16, 4: 6.18, 5: 6.21, 6: 6.24, 7: 6.29, 8: 6.30, 9: 6.34, 10: 6.35.\par
\noindent PS8: 1: 6.33, 2: 6.39, 3: 6.40, 4: 6.41, 5: 7.2, 6: 7.3, 7: 7.4.\par
\noindent PS9: 1: 7.5, 2: 7.7, 3: 7.9, 4: 7.11, 5: 8.1, 6: 8.3, 7: 8.4, 8: 8.5, 9: 8.6.\par
\noindent PS10: 1: 8.8, 2: 8.10, 3: 8.11, 4: 8.12, 5: 9.1, 6: 9.3, 7: 9.4, 8: 9.6, 9: 9.12.\par
PS1-8 (K\&K 1.20): 径向与切向加速度分别指 $\uv r,\uv\theta$ 分量. PS2-1 (2.1): 质量从静止出发. PS2-9 (2.17): 原卷勘误为 (b) $\tan\theta>\mu$, (c) $\tan\theta<\mu$. PS3 没有完整自编题, 仅索引中的九项. 不能因为它们引用同一本书而排除其他卷的独立 MIT 题.

\originalsection{作业 1}
\begin{exercise}
原题 9. 相距 $d$ 两城同轨两列车同时以速率 $v$ 相向出发. 蜜蜂从 A 城车头以 $u>v$ 飞向 B 车, 遇车立即反向, 不断往返直到两车相撞. (a) 第一次飞行记 $d_0$, 随后为 $d_1,d_2,\ldots$, 求每段距离并求和. (b) 用很少计算的另一方法求总路程.
\end{exercise}
\begin{solution}
(a) 第一次遇车时间 $d/(u+v)$, 路程 $d_0=ud/(u+v)$. 每次遇车后两车间距乘 $q=(u-v)/(u+v)$, 故 $d_n=d_0q^n$. 求和 $D=d_0/(1-q)=ud/(2v)$. (b) 两车相撞时间 $d/(2v)$, 蜂一直以速率 $u$ 飞, 因而 $D=u d/(2v)$. 无数次瞬时反向是理想化, 蜂的加速度未被建模.
\end{solution}
\begin{exercise}
原题 10. 量纲分析把未知量写成相关变量与常数的幂乘积. (a) 求引力不稳定驱动的质量 $M$、半径 $R$ 恒星振动频率尺度. (b) 球 $M,R$ 以 $v$ 在密度 $\rho$ 介质中运动, 求阻力尺度. (c) 地表球受 $F=bv^2$, 求终端速度并给另一推导. (d) 单摆 $M,L,\theta_0$ 的频率尺度, 能否确定对无量纲角的依赖?
\end{exercise}
\begin{solution}
(a) $f\sim\sqrt{GM/R^3}$, 无法由量纲确定结构相关系数. (b) 在惯性阻力、只取 $\rho,R,v$ 时 $F\sim\rho R^2v^2$; 若把 $M$ 也允许进入, 无量纲比 $M/(\rho R^3)$ 可进入未知函数, 单靠量纲不能唯一排除它, 更不能确定阻力系数. (c) $v_t\sim\sqrt{Mg/b}$, 力平衡 $Mg=bv_t^2$ 在该模型中给精确 $v_t=\sqrt{Mg/b}$. (d) $f=\sqrt{g/L}\,f_1(\theta_0)$, 量纲不能决定 $f_1$, 小角微分方程才给 $1/(2\pi)$.
\end{solution}
\originalsection{作业 2}
\begin{exercise}
原题 11. 三个等边三角形拼成七根无质量铰接杆桁架. 下层节点 A,B,C 在 $x=0,l,2l$, 上层 D,E 在 $x=l/2,3l/2$, 高 $\sqrt3l/2$. 杆为 AD,DB,BE,EC,AB,BC,DE. 质量 $M$ 在中央 B, 两端 A,C 只给竖直支持. (a) 求各杆两端力大小、方向及拉压. (b) 加分: 推广至 $n$ 个三角形、$2n+1$ 根杆.
\end{exercise}
\begin{solution}
(a) 对称支承各 $Mg/2$. 二力杆的端力均沿杆. A 的竖直平衡给 AD 压力 $Mg/\sqrt3$, 水平给 AB 拉力 $Mg/(2\sqrt3)$. B 的平衡给 DB、BE 拉力各 $Mg/\sqrt3$. D 的平衡给 DE 压力 $Mg/\sqrt3$. 右半对称: EC 压力 $Mg/\sqrt3$, BC 拉力 $Mg/(2\sqrt3)$. 每杆对节点的力, 拉杆朝另一端, 压杆背离另一端; 杆所受节点力方向相反.

(b) 原加分问未给推广图, 这里明确采用与 (a) 相同的交替三角链、奇数 $n$, 两端同高, 中央节点承载. 编节点 $j=0,\ldots,n+1$, 坐标 $(jl/2,0)$ (偶数) 或 $(jl/2,\sqrt3l/2)$ (奇数), 中央 $k=(n+1)/2$. 对角杆 $j\to j+1$ 的有向轴力 (拉正压负) 为
\begin{equation}
D_j=\frac{Mg}{\sqrt3}\begin{cases}(-1)^{j+1},&j<k,\\(-1)^j,&j\ge k.\end{cases}
\end{equation}
水平杆 $j\to j+2$ 为 $C_j=(-1)^jMg\min(j+1,n-j)/(2\sqrt3)$, $j=0,\ldots,n-1$. 用各节点两方向平衡可逐次推出这两个式, 也可直接代入, 中央两对角竖直分量相加恰承托 $Mg$. 偶数 $n$ 的端点不同高且几何中心未必为节点, 需指定载荷节点或横梁才有唯一的对应推广, 不把上述另加条件隐作原题.
\end{solution}
\begin{exercise}
原题 12. 尽量不用计算器作数量级估算. (a) 全球人口伸臂牵手可绕地球几圈, 地球半径约 $6\times10^6$ m? (b) 星体每年 1 角秒的角位移投影到距 $4\times10^8$ m 月面, 线速度像蜗牛、跑步者还是喷气机? (c) 推出汽车初速、变形区长度与减速度尺度, 30 mph 撞击限 20$g$ 至少需多长, 与车长比? (d) 最近恒星约距 4 光年, $c=3\times10^8$ m/s, 典型恒星 $10^{30}$ kg, 估附近恒星质量密度并与空气比较.
\end{exercise}
\begin{solution}
(a) 明确用 2008 年估值人口 $6\times10^9$, 每人臂展约 2 m, 链长 $1.2\times10^{10}$ m, 周长約 $3.8\times10^7$ m, 约 300 圈. 非当前人口统计. (b) $1''=\pi/(180\times3600)\simeq5\times10^{-6}$ rad, 一年约 $3\times10^7$ s, 所以速度约 $7\times10^{-5}$ m/s, 更接近蜗牛. (c) 量纲 $a\sim v^2/d$, 匀减速精确 $d=v^2/(2a)$, 30 mph 约 13 m/s, $a=20g\simeq200$ m/s$^2$, 得 $d\simeq0.45$ m, 小于几米车长但需真实人体和结构模型才能评估. (d) 4 光年约 $4\times10^{16}$ m, 每个此尺度立方体一个星, 密度约 $10^{30}/(4\times10^{16})^3\sim10^{-20}$ kg/m$^3$, 比约 1 kg/m$^3$ 空气低 20 数量级, 对应粗略间距模型.
\end{solution}
\originalsection{作业 4}
\begin{exercise}
原题 9. $n$ 个球从底向上编号, $m_1\gg m_2\gg\cdots\gg m_n$, 最重球在底部, 从离地 $H=1$ m 静止释放. 忽略空气阻力和真实弹性极限, 各次碰撞完全弹性. (a) 顶球弹多高, 达 1 km 需几球? (b) 达地球逃逸速率需几球, 可用量纲估算逃逸速度.
\end{exercise}
\begin{solution}
落地前共同向下速率 $u=\sqrt{2gH}$. 底球弹起速率 $u$, 逐级视大球为移动墙, 有 $v_1=u$, $v_j=2v_{j-1}+u=(2^j-1)u$. 因而 (a) 顶球高度约 $(2^n-1)^2H$, $n=5$ 得 961 m, $n=6$ 得 3969 m, 所需 6 球. 忽略球半径与原顶球高度, 且要求碰撞次序分离. (b) $v_e\sim\sqrt{GM_E/R_E}$, 精确球对称模型为 $\sqrt{2gR_E}\simeq1.1\times10^4$ m/s. 需 $2^n-1\ge v_e/u\simeq2450$, 故 12 球. 高速或高弹跳时恒定 $g$ 高度式失效, 此处逃逸判断直接比较地表速度.
\end{solution}
\begin{exercise}
原题 10. 均匀链 $M,L,\lambda=M/L$, 长直段初竖直位于桌边上方, 下端经 U 形折返接在固定桌边. 松开上端, 链只作竖直运动, 忽略摆动与反弹, 最终垂挂于桌下. 求桌支承力随时间. 原图并非链落在桌面堆积.
\end{exercise}
\begin{solution}
题目未给折返与末端停止的耗散模型, 故仅凭 $M,L,g$ 不能宣称一个普遍的唯一力--时间曲线. 以下给出明确的连续柔性、折返无耗散模型. 折返点向下移 $y$, 自由端降 $2y$, 静止支段长 $y$, 运动段 $L-y$, 速率 $v=2\dot y$. 势能下降为 $\lambda g(2Ly-y^2)$, 动能为 $\lambda(L-y)v^2/2$, 因而
\begin{equation}
v^2=\frac{4gy(L-y/2)}{L-y},\qquad
 t(y)=\int_0^y\sqrt{\frac{L-s}{g s(L-s/2)}}\dd s.
\end{equation}
折返处相对通量为 $\lambda v/2$, 速度由向下 $v$ 到零, 两側张力合为 $\lambda v^2/2$, 每侧 $T=\lambda v^2/4$. 静止段顶部支力
\begin{equation}
F_S(t(y))=\lambda gy+\frac{\lambda v^2}{4}
=\lambda g\left[y+\frac{y(L-y/2)}{L-y}\right],\qquad0<y<L.
\end{equation}
$y\to L$ 时该理想连续模型的速度与力奇异, 不能给出有限末端峰值. 全链静止后支力为 $Mg$, 中间停止冲击须加弹性与耗散模型. 随附数值积分仅复现上述已声明模型, 不修复原题边界欠规定.
\end{solution}
\originalsection{作业 5 与作业 6}
\begin{exercise}
PS5 原题 10. 估算太阳能能否满足波士顿需求. (a) 人口约 600 万, 估用电平均功率. (b) 太阳光度 $4\times10^{26}$ W、日地距 $1.5\times10^{11}$ m, 都市面积 11000 km$^2$, 电池效率约 10\%, 计日夜、倾角、云及大气损失估可用通量. (c) 求供全部用电所需面积比例及障碍. (d) 对宿舍与四人独栋住宅重算, 屋顶够否, 是否每栋屋顶均能供全楼? (e) 原题引 WRI 美国人均总能源 80000 kWh/年, 换成功率并解释与用电差别, 总能源需求能否覆盖?
\end{exercise}
\begin{solution}
明确选每人平均电功率 1 kW 的估算, (a) $P_e\sim6$ GW. (b) 入射 $F_\odot=L_\odot/(4\pi d^2)\simeq1415$ W/m$^2$. 取日夜因子 $1/2$, 几何 $1/2$, 云与大气 $1/2$, 再乘效率 0.1, 年平均约 18 W/m$^2$. 这些因子是估算, 不是实测当地辐照. (c) 面积約 $3.4\times10^8$ m$^2$=340 km$^2$, 约 3\%; 还需储能、输电、季节平衡、可用土地等. (d) 示例宿舍 200 人、屋顶 1000 m$^2$, 需 200 kW 而可供約 18 kW; 四人家需 4 kW, 需约 220 m$^2$ 有效面板, 100 m$^2$ 屋顶不够. 真实建筑需给人数、耗电和屋顶面积, 不能凭“宿舍”给唯一数字; 层数越多屋顶与人口比一般越小. (e) $80000\times1000/(365\times24)\simeq9130$ W/人, 包含交通、供热和工业等非电能源. 都市需約 55 GW, 同估算通量需約 3100 km$^2$, 约 28\%. 这是题设历史能源量的几何估算, 不等于当代可实施能源规划.
\end{solution}
\begin{exercise}
PS6 原题 7. 圆柱质量 $M$、半径 $R$, 靠竖直粗糙墙静止, 一根从墙高处到圆柱右上方的绳与圆柱相切, 绳与竖直墙夹角 $\theta$. 墙摩擦系数 $\mu$. 求最小 $\theta$, 并解释为何为最小而非最大, 圆柱提高或降低会如何.
\end{exercise}
\begin{solution}
设绳张力 $T$, 墙法向力 $N$, 墙摩擦 $f$ 向上. 绳对圆心矩为 $RT$, 墙摩擦矩反向为 $Rf$, 故 $f=T$. 水平 $N=T\sin\theta$, 竖直 $T\cos\theta+f=Mg$, 得 $T=Mg/(1+\cos\theta)$. 静摩擦要求 $T\le\mu T\sin\theta$, 即 $\sin\theta\ge1/\mu$. 在图示 $0<\theta\le\pi/2$ 范围, 需 $\mu\ge1$, $\theta_{\min}=\arcsin(1/\mu)$. 圆柱相对固定上绳点升高使绳更倾斜、$\theta$ 增大, 法向压紧相对所需摩擦更足; 降低使角减小而可失稳. 若绳作用线过圆心, 会变成不同题目.
\end{solution}
\originalsection{作业 7 与作业 8}
\begin{exercise}
PS7 原题 11. 均匀实心弹力球 $M,R$, 初水平速度 $v_x$, 自转轴垂直轨迹平面. 在地面作完全弹性粗糙反弹, 竖直速率大小不变. (a) 求一次反弹后水平速度与自转. (b) 求第二次后值并描述多次行为; 可自行实验核对但本稿不宣称已做实物实验.
\end{exercise}
\begin{solution}
取向右为 $v$, 逆时针为 $\omega$, $I=\kappa MR^2$, $\kappa=2/5$. 将“无滑弹性反弹”明确解释为接触切向相对速度反号的粗糙弹性冲量模型, $u'=-(v+R\omega)$. 切向冲量 $J=M(v'-v)$, 角冲量 $RJ=I(\omega'-\omega)$, 解得
\begin{equation}
v'=\frac{3v-4R\omega}{7},\qquad
\omega'=\frac{-10v/R-3\omega}{7}.
\end{equation}
该线性变换平方为单位, 所以 (b) 第二次恢复原 $v,\omega$, 理想多次反弹交替两状态. 如果“无滑”被理解为碰后永久粘着 $u'=0$, 切向动能会损失, 不能同时满足完全弹性; 原词语存在此模型区别.
\end{solution}
\begin{exercise}
PS8 原题 8. 均匀矩形薄片水平边 $a$, 竖直边 $b$, 初自由静止, 左上角受 $F\uv z$, 右下角受 $-F\uv z$ 持续短时 $\Delta t$. 为写惯量定义其质量为 $M$. (a) 求质心 $x,y$ 轴惯量并核垂直轴定理. (b) 求末角速度和角动量, 证明 $a\ne b$ 时不平行, 角速度沿另一条对角线.
\end{exercise}
\begin{solution}
(a) $I_x=Mb^2/12$, $I_y=Ma^2/12$, $I_z=M(a^2+b^2)/12=I_x+I_y$. (b) 两冲量合为零, 角冲量为 $\vct L=F\Delta t(b\uv x+a\uv y)$. 主轴关系给 $\vct\omega=(12F\Delta t/M)(\uv x/b+\uv y/a)$, 方向等于 $a\uv x+b\uv y$, 即从左下到右上的另一对角线. 两向量的分量比互换, 仅当 $a=b$ 才平行.
\end{solution}
\begin{exercise}
PS8 原题 9. 均匀圆盘在桌上无滑摇摆, 与水平夹角 $\alpha$, 可把 $\alpha$ 和质心位置视为常数. 自旋 $\omega_s$, 接触点沿桌面转动角频率 $\Omega$. (a) 证总角速度大小 $\Omega\sin\alpha$, 沿接触点指向质心. (b) 证 $\Omega=2\sqrt{g/(R\sin\alpha)}$ 并讨论极限. (c) 证俯视盘面的表观转速为 $(1-\cos\alpha)\Omega$. 原图半径记 $R$, 原公式误用未定义 $L$, 此处统一用 $R$.
\end{exercise}
\begin{solution}
(a) 接触点与质心瞬时速度都为零, 所以瞬时转轴沿连线. 写竖直进动加盘法向自旋, 沿盘法向投影无分量给 $\omega_s=-\Omega\cos\alpha$, 剩余连线分量为 $\Omega\sin\alpha$. (b) 连线是盘内直径主轴, $I=MR^2/4$, 角动量 $L=I\Omega\sin\alpha$. 第\ref{ch:angular}章力矩比较得 $\Omega^2=4g/(R\sin\alpha)$. $\alpha\to0$ 模型发散, $\alpha\to\pi/2$ 极限 $2\sqrt{g/R}$; 恰直立时力矩比较同时出现零因子, 应理解为极限而非唯一进动值. (c) 随接触方向旋转的几何方位角增长率为 $\Omega$, 盘面图案相对该方向自旋为 $-\Omega\cos\alpha$, 合为 $(1-\cos\alpha)\Omega$.
\end{solution}
\originalsection{作业 9 与作业 10}
\begin{exercise}
PS9 原题 10. 水分子两个氢各 $m_H=1$ amu, 氧质量 $16m_H$, 键长 $r=9.6\times10^{-11}$ m, 夹角 $104.5^\circ$, 1 amu=$1.66\times10^{-27}$ kg. 各原子视点质量. (a) 求质心与过质心主轴, 计算主轴惯量张量. (b) 哪轴无矩转动不稳定?
\end{exercise}
\begin{solution}
令氧在原点, $x$ 沿角平分线朝两氢, 氢坐标 $(r c,\pm r s,0)$, $c=\cos52.25^\circ,s=\sin52.25^\circ$. 质心为 $(rc/9,0,0)$. 对称给各非对角元零, 三轴为主轴. (a) $I_x=2m_Hr^2s^2$, $I_y=(16/9)m_Hr^2c^2$, $I_z=I_x+I_y$. (b) 数值有 $I_y<I_x<I_z$, 中间惯量轴 $x$ 不稳定, 理由为欧拉方程的小扰动指数增长, 见第\ref{ch:angular}章.
\end{solution}
\begin{exercise}
PS10 原题 10. (a) 推导角速度随时间变的旋转系中额外惯性力. (b) 理想花滑者双手各质量 $M$, 其余质量忽略, 初半径 $R$, 自旋向上 $\omega$. 双手缓慢向内, 径向速率 $\dot r$, 在共转系中用全部惯性力求 $\dot\omega$. (c) 从惯性系角动量导同式. (d) 半径收至一半, 角速率变多少?
\end{exercise}
\begin{solution}
(a) 第\ref{ch:frames}章导数给欧拉力 $-M\dot{\vct\omega}\times\vct r$. (b) 共转手无相对角运动, 切向真实力在理想径向拉臂模型为零, 科里奥利分量 $-2M\omega\dot r$, 欧拉分量 $-Mr\dot\omega$, 相加为零. 故 $\dot\omega=-2\omega\dot r/r$, 初时 $r=R$. 离心力只在径向, 不进入此切向式. (c) $L=2Mr^2\omega$ 恒定, 求导为同式. (d) $\omega r^2$ 恒定, 末角速率为四倍.
\end{solution}
\originalsection{开放项目与实际数值示例}
原项目占 5\%, 最多四人组, 可以构建机械装置、实验测量或数值模拟, 要检验课堂原理而非只做展示. 原要求第 14 讲提交不超过一页、12 pt、单倍行距的提案, 包含题名、成员、目标、设计、与力学关系、所需资源和获取方法; 先经教师接受. 第 21 讲交 3--5 页报告与图片、视频或模拟等材料, 按物理解释、设计执行、分析及创造性评分. PS6 另给提案日期 2008-10-29、报告日期 2008-11-26. 这些是当年课程规定, 不向读者虚构当前需要的 MIT 批准.

本册附一个实际运行的数值项目: 长度 1 m 单摆初角 0.5 rad、初角速零, 对 $\ddot\theta=-g\sin\theta$ 使用四阶 Runge--Kutta, 不采用小角近似. 在第一次过零附近线性插值, 四倍过零时间估周期, 同时记录离散能量相对漂移. 脚本 \texttt{tools/numerics.py} 和输出 CSV 纳入源码. 不同时间步的收敛仅验证该初值数值复现, 不证明一般周期公式. 与椭圆积分级数和小角公式的比较在运行表中给出, 没有伪造实物测量.
\begin{center}\begin{tabular}{rrr}\toprule
步长 / s & 单摆周期 / s & 10 s 最大相对能量漂移 \\\midrule
0.0200 & 2.03786885 & $3.891\times10^{-7}$ \\
0.0100 & 2.03786764 & $1.240\times10^{-8}$ \\
0.0050 & 2.03786785 & $4.034\times10^{-10}$ \\
0.0025 & 2.03786791 & $1.360\times10^{-11}$ \\
\bottomrule\end{tabular}\end{center}
独立 Simpson 积分有限振幅周期为 2.03786792 s, 小角周期为 2.00606668 s. 周期插值本身只有二阶, 不能把过零时间误差也称为四阶.
链条采用前述无耗散折返模型, $L=1$ m. 用 $y=L\sin^2\phi$ 消去端点根号, 实际积分得如下支力曲线点; 不把末端奇异力当成实测峰值.
\begin{center}\begin{tabular}{rrr}\toprule
$y/L$ & $t$ / s & $F_S/(Mg)$ \\\midrule
0.00 & 0.000000 & 0.0000 \\
0.10 & 0.200179 & 0.2056 \\
0.25 & 0.311917 & 0.5417 \\
0.50 & 0.428056 & 1.2500 \\
0.75 & 0.502842 & 2.6250 \\
0.90 & 0.531449 & 5.8500 \\
0.99 & 0.540688 & 50.9850 \\
\bottomrule\end{tabular}\end{center}
折返完成时刻的有限极限为 0.54098864 s; 此时理想支力发散, 最终静止的 $Mg$ 不是该曲线连续延拓.
另运行期末第 7 题的中心力初值: $A=m=L=1$, 初 $r=1.01$, 纯切向速率 $1/1.01$. 步长 0.001, 积分 20 时间单位. 实得径向周期 2.56495053, 线性预测 $2\pi/\sqrt6=2.56509966$; 有限振幅差约 5.81e-05. 最大相对能量漂移 1.03e-14. 数据仅验证这一接近圆轨道的复现, 一般结论由正文推导保证.

